\section{Methodology}
\label{sec:methodology}
%
In the following, we will use the ideas of \textcite{deliot2023} to derive a fast, stable, and visually convincing glint renderer, that converges to any given configuration of the Cook-Torrance model.


\subsection{Shading Model}
\label{sec:shading}
%
Similar to \textcite{zirr2016}, \textcite{deliot2023} model the surface as consisting of mesoscopic details covered with microscopic details or \emph{microfacets}.
% TODO: Add illustrating figure
%
Subsequently, we will denote these larger-scale features as \emph{meso-facets}, and 
% we will refer to % NOTE: May flow better
meso-facets that contribute to the reflectance of light as \emph{glints}. 


% ?: Use we or \textcite{deliot2023}
Therefore, \textcite{deliot2023} start from the Rendering Equation \parencite{kajiya1986, immel1986}:
%
\begin{equation}
    L_o(\vec{\omega_o}) =
    \int_{\Omega}
        \rho_P(\vec{\omega_i}, \vec{\omega_o})~
        L_i(\vec{\omega_i})~
        \dotproduct{\vec{\omega_i}}{\vec{n}}~
    d\vec{\omega_i}.
\end{equation}
%
% NOTE: Maybe write out BRDF
The outgoing radiance $L_o$ emitted in direction $\vec{\omega_o}$ is the sum over the hemisphere $\Omega$ of incoming radiance $L_i$ modulated by the BRDF $\rho_P$, whereas $P$ being the pixel footprint and $\vec{n}$ being the surface normal.
%
$\rho_P$ is based on a Cook-Torrance model \parencite{cook1982} consisting of the Fresnel reflectance term $F$, the geometrical attenuation term $G$, and the distribution term $D$:
%
% NOTE: Maybe clarify that this is the specular component
\begin{equation}
    \rho_P(\vec{\omega_i}, \vec{\omega_o}) = \frac{
        F(\dotproduct{\vec{\omega_i}}{\vec{\omega_h}})~
        G(\vec{\omega_i}, \vec{\omega_o})~
        D_P(\vec{\omega_h})
    }{
        4~
        \dotproduct{\vec{\omega_i}}{\vec{n}}~
        \dotproduct{\vec{\omega_o}}{\vec{n}}
    },
\end{equation}
%
% TODO: Rewrite this
with the half-vector $\vec{\omega_h}=\flatnormalize{\vec{\omega_i} + \vec{\omega_o}}$ describing the surface normal at which perfect reflection would occur and $D_P(\vec{\omega_h})$ as the probability for a microfacet on a meso-facet inside the pixel footprint to be aligned along this half-vector $\vec{\omega_h}$ and thereby reflect light.
%
This differs from the original distribution term $D(\vec{\omega_h})$ which only represents the probability of a normally distributed microfacet being aligned that way. We get
%
\begin{equation}
\begin{aligned}
    D_P(\vec{\omega_h})
    % TODO: Rename probabilities / explain probabilities
    & = && \probability{\mathrm{glint}} && \cdot && \probabilitygiven{\mathrm{micro}}{\mathrm{glint}} \\
    & = && \frac{b(N(P), p)}{N(P)} && \cdot && D(\vec{n}),
\end{aligned}
\end{equation}
%
where $b(n, p)$ is a random sampling of the binomial distribution with $n$ trials and a success probability of $p$. $N(P)$ is the number of meso-facets inside the pixel footprint and therefore $b(N(P), p)$ counts the number of glints inside $P$.

We then choose the glint probability $p = \frac{D(\vec{\omega_h})}{D(\vec{n})}$ to ensure $\expectation{D_P(\vec{\omega_h})} = D(\vec{\omega_h})$:
%
\begin{equation}
    % TODO: Rename to D’_P
    \expectation{D_P(\vec{\omega_h})}
    % = \frac{\E{b(N(P), p)}}{N(P)} D(\vec{n})
    = \frac{N(P) \, p}{N(P)} D(\vec{n})
    = \frac{D(\vec{\omega_h})}{D(\vec{n})} D(\vec{n})
    = D(\vec{\omega_h}).
\end{equation}
%
Furthermore, \textcite{deliot2023} introduce an artistic control parameter $R \in (0,1]$ that scales the glint probability. This leaves us with:
%
\begin{equation}
\label{eq:glint-distribution-term}
    D_P(\vec{\omega_h})=\frac{b(N(P), p)}{N(P)\,R} \cdot D(\vec{n}),
    ~\text{with}~
    p = R~\frac{D(\vec{\omega_h})}{D(\vec{n})}
\end{equation}


$R$ can not directly be linked to a physical quantity but yet can be interpreted as a measure of \emph{meso-facet roughness}: Lowering $R$ results in fewer but brighter glints per meso-facets and thus visually makes the glints less stable.


\subsubsection{Convergence}
\label{sec:convergence-proof}

We observe that for large meso-facet count $N(P)$
%
\begin{equation}
    \lim_{N(P) \to \infty} D_P(\vec{\omega_h})
    = \frac{D(\vec{n})}{R} \lim_{N(P) \to \infty} \frac{b(N(P), p)}{N(P)}
    = D(\vec{\omega_h}),
\end{equation}
%
because, according to the law of large numbers, as $n$ approaches infinity, $\lim_{n \to \infty}\frac{b(n,p)}{n}=\expectation{\frac{b(n,p)}{n}}=p$.


Hence, we have shown that our shading model converges to the chosen base Cook-Torrance microfacet model, and we thereby satisfy the \ref{req:convergence} requirement.